\section{Benchmark}

\benchmarkname{} contains short single-turn prompts. An item is eligible when
it satisfies these acceptance criteria:

\begin{enumerate}
  \item the prompt can be understood without external context;
  \item the expected answer is objective and unambiguous;
  \item the answer can be derived without tools;
  \item a deterministic scorer can check the response; and
  \item failure is visible in a public-facing AI product.
\end{enumerate}

The current paper candidate split, \paperSplit{}, is seeded from an 80-item
hard-obvious panel with 10 candidate items from each of eight families. The
paper-readiness gate checks the manifest-scoped item cards, scorer-gold
fixtures, and split structure before these items support claims. The table
below is generated from the manifest used by the paper build.

\begin{table}[t]
  \centering
  \caption{Current \paperSplit{} composition generated from the paper
  manifest.}
  \label{tab:dataset-composition}
  % Generated by scripts/build_paper_assets.py. Do not hand-edit.
\begin{tabular}{llr}
\toprule
Family & Subfamily & Items \\
\midrule
Arithmetic & Numeric comparison & 3 \\
Arithmetic & Small integer arithmetic & 2 \\
Arithmetic & Unit conversion and small calc & 5 \\
Char count & Single-letter count & 10 \\
Constraints & Object must be present & 10 \\
Format & Exact JSON schema & 5 \\
Format & JSON field & 5 \\
Negation & Not-choice & 5 \\
Negation & Without constraint & 5 \\
Ordering & Alphabetical sort & 5 \\
Ordering & Numeric sort & 5 \\
Spelling & Remove letter & 9 \\
Spelling & Replace letter & 1 \\
Word count & Comma-list count & 10 \\
\midrule
\multicolumn{2}{r}{Total} & 80 \\
\bottomrule
\end{tabular}

\end{table}

\paragraph{Task families.}
The task families are character counting, spelling transforms, small
arithmetic or numeric comparison, word and list counting, ordering, negation,
format compliance, and constraint awareness. Each family must have reviewed
item cards and scorer-gold evidence before supporting a final claim.

\paragraph{What it does not measure.}
\benchmarkname{} does not measure general intelligence, factual knowledge,
safety, tool use, long-context behavior, retrieval quality, multi-turn agent
behavior, or expert reasoning. It is a public-facing reliability check, not a
comprehensive capability exam.
